La siguiente tabla indica el valor de las variables de un problema de transporte con las ofertas de los orígenes $O_i$ y las demandas de los destinos $D_j$. Indica cuáles son los vectores de tasas de sustitución de las variables $x_{14}$ y $x_{22}$.

Nota: la celda correspondiente a la fila $i$ y a la columna $j$ contiene el valor de la variable $x_{ij}$, e indica la cantidad de producto que se transporta desde el origen $i$ al destino $j$.

\begin{table}[h]
\begin{center}

\begin{tabular}{ccccccc}
                      &                         &                        &                        &                        &                         & $O_i$ \\ \cline{2-6}
\multicolumn{1}{l|}{} & \multicolumn{1}{l|}{10} & \multicolumn{1}{l|}{}  & \multicolumn{1}{l|}{}  & \multicolumn{1}{l|}{}  & \multicolumn{1}{l|}{}   & 10    \\ \cline{2-6}
\multicolumn{1}{l|}{} & \multicolumn{1}{l|}{5}  & \multicolumn{1}{l|}{5} & \multicolumn{1}{l|}{5} & \multicolumn{1}{l|}{}  & \multicolumn{1}{l|}{}   & 15    \\ \cline{2-6}
\multicolumn{1}{l|}{} & \multicolumn{1}{l|}{}   & \multicolumn{1}{l|}{}  & \multicolumn{1}{l|}{5} & \multicolumn{1}{l|}{5} & \multicolumn{1}{l|}{10} & 20    \\ \cline{2-6}
$D_j$                 & 15                      & 5                      & 10                     & 5                      & 10                      &      
\end{tabular}
    
\end{center}
\end{table}